% Part: computability
% Chapter: computability-theory
% Section: universal-part-function

\documentclass[../../../include/open-logic-section]{subfiles}

\begin{document}

\olfileid[tr]{cmp}{thy}{uni}
\olsection{Evrensel Kısmi Hesaplanabilir Fonksiyon}

\begin{thm}
\ollabel{thm:univ-comp}
Evrensel bir kısmi hesaplanabilir $\fn{Un}(e,x)$ fonksiyonu vardır. Başka bir
deyişle aşağıdaki özellikleri taşıyan bir $\fn{Un}(e,x)$ fonksiyonu vardır:
\begin{enumerate}
\item $\fn{Un}(e,x)$ kısmi hesaplanabilirdir.
\item $f(x)$ herhangi bir kısmi hesaplanabilir fonksiyonsa, her $x$ için
  $f(x)\simeq\fn{Un}(e,x)$ olacak bir $e$ doğal sayısı vardır.
\end{enumerate}
\end{thm}

\begin{proof}
$U$ ile $T$, Kleene normal biçim teoremindeki gibi olmak üzere
(\olref[nfm]{thm:normal-form}),
$\fn{Un}(e,x)\simeq U(\umin{s}{T(e,x,s)})$ diyelim.
\end{proof}

\begin{explain}
Bu, kısmi hesaplanabilir fonksiyonların etkin bir numaralandırmasına sahip
olduğumuzu kesin biçimde söylemenin yoludur. Düşünce şudur:
$f_e(x)=\fn{Un}(e,x)$ ile tanımlanan fonksiyona $f_e$ dersek
$f_0,f_1,f_2,\dots$ dizisi bütün kısmi hesaplanabilir fonksiyonları içerir;
üstelik $f_e(x)$, $e$ ile $x$ bakımından ``birbiçimli'' olarak
hesaplanabilir. Kolaylık için tek yerli fonksiyonlar bakımından evrensel olan
ikili bir fonksiyon kullanıyoruz; fakat sayı dizilerini kodlayarak bunu daha
çok argümana kolayca genelleyebiliriz. Örneğin $f(x,y,z)$ üç yerli kısmi
özyinelemeli bir fonksiyonsa
$g(x)\simeq f((x)_0,(x)_1,(x)_2)$ tek yerli kısmi özyinelemeli bir
fonksiyondur.
\end{explain}

\end{document}

